\documentclass[aspectratio=169,12pt]{beamer}

\usepackage[T2A]{fontenc}
\usepackage[utf8]{inputenc}
\usepackage[russian]{babel}
\usepackage{graphicx}
\usepackage{xcolor}
\usepackage{array}

\definecolor{mintbg}{RGB}{208,229,205}
\definecolor{mintline}{RGB}{148,180,144}
\definecolor{textgray}{RGB}{60,60,60}

\setbeamertemplate{navigation symbols}{}
\setbeamercolor{normal text}{fg=black,bg=white}
\setbeamercolor{frametitle}{fg=black,bg=mintbg}
\setbeamerfont{frametitle}{size=\Large}
\setbeamersize{text margin left=8mm,text margin right=8mm}
\newcommand{\MainSlidesPdf}{output/presentation.pdf}

\newcommand{\NotesFrame}[3]{%
\begin{frame}[t]{#1}
    \centering
    \fbox{\includegraphics[page=#2,width=0.75\linewidth,height=0.34\textheight,keepaspectratio]{\MainSlidesPdf}}

    \vspace{0.4em}
    \raggedright
    \small
    #3
\end{frame}
}

\title{Текст выступления}
\subtitle{Программная система фиксации городских событий}
\author{Берёзкина Е.\,В.}
\date{}

\begin{document}

\begin{frame}
    \titlepage
\end{frame}

\NotesFrame{Слайд 1. Титульный слайд}{1}{
Добрый день! Тема моей работы --- <<Разработка программной системы фиксации городских событий>>.

\vspace{0.4em}

Работа выполнена в рамках преддипломной практики и объединяет результаты двух семестров: серверную часть, веб-панель диспетчера и мобильное приложение для жителей.
}

\NotesFrame{Слайд 2. Анализ текущей проблемы}{2}{
В настоящее время в городе Кирове отсутствует единый цифровой канал для сообщения о проблемах городской инфраструктуры. Информация поступает по разрозненным каналам: телефон, почта, личные сообщения.

\vspace{0.4em}

Это приводит к потере обращений, затруднённой маршрутизации и отсутствию обратной связи. Без карты диспетчер не видит общую картину, а у жителей нет удобного инструмента для фиксации проблем на месте.
}

\NotesFrame{Слайд 3. Цель и задачи работы}{3}{
Цель работы --- разработать программную систему для фиксации и обработки обращений о проблемах городской инфраструктуры города Кирова.

\vspace{0.4em}

Система включает три компонента: серверную часть с базой данных, диспетчерскую веб-панель с интерактивной картой и кроссплатформенное мобильное приложение для жителей.

\vspace{0.4em}

Задачи: анализ предметной области, проектирование структуры и базы данных, реализация сервера, веб-панели и мобильного приложения, тестирование.
}

\NotesFrame{Слайд 4. Анализ существующих решений}{4}{
Были рассмотрены существующие аналоги: <<Госуслуги. Решаем вместе>>, <<Наш город>> (Москва), <<Добродел>> и FixMyStreet.

\vspace{0.4em}

Ни одно из них не обеспечивает полный набор функций для Кирова: автоматическое определение геолокации, карту обращений, отслеживание статуса и механизм дедупликации. Поэтому было принято решение разрабатывать собственную систему.
}

\NotesFrame{Слайд 5. Технологический стек}{5}{
Серверная часть --- ASP.NET Core на C\# с PostgreSQL и расширением PostGIS для геопространственных запросов.

\vspace{0.4em}

Веб-панель --- React с TypeScript и Яндекс.Картами для интерактивной карты обращений.

\vspace{0.4em}

Мобильное приложение --- React Native с Expo. Единая кодовая база для Android и iOS, преемственность подходов с веб-панелью.
}

\NotesFrame{Слайд 6. Структура системы}{6}{
Вот общая картина: мобильное приложение и веб-панель --- это два клиента, которые обращаются к одному серверу на C\#.

\vspace{0.4em}

Сервер работает с PostgreSQL и PostGIS. Оба клиента используют единый REST API, что обеспечивает согласованность данных.
}

\NotesFrame{Слайд 7. Особенности проектирования базы данных}{7}{
Основной сущностью базы данных является обращение: описание, координаты, статус, дата.

\vspace{0.4em}

Обращение связано с категорией проблемы, пользователем, медиафайлами, историей статусов и ответственными службами. Такая структура позволяет полноценно сопровождать заявку на всех этапах обработки.
}

\NotesFrame{Слайд 8. Реализация: дашборд диспетчера}{8}{
Дашборд содержит сводные показатели по обращениям, графики распределения по категориям и статусам, а также список последних поступивших записей.

\vspace{0.4em}

Дашборд выполняет функцию стартовой аналитической панели для диспетчера.
}

\NotesFrame{Слайд 9. Реализация: интерактивная карта диспетчера}{9}{
Центральным рабочим экраном диспетчера является интерактивная карта обращений на базе Яндекс.Карт.

\vspace{0.4em}

На карте отображаются маркеры обращений, цвет зависит от статуса. Диспетчер может фильтровать по состояниям и быстро находить проблемные точки.
}

\NotesFrame{Слайд 10. Реализация: список обращений}{10}{
Страница содержит табличный перечень записей с основными полями, вкладки по статусам и средства быстрого поиска.

\vspace{0.4em}

Это основной режим работы со структурированным перечнем поступивших сообщений.
}

\NotesFrame{Слайд 11. Реализация: карточка обращения}{11}{
Карточка содержит полную информацию по заявке: описание, категория, статус, история обработки, мини-карта и данные ответственной службы.

\vspace{0.4em}

Также диспетчеру доступна форма смены статуса с добавлением комментария. Этот экран обеспечивает непосредственную диспетчерскую обработку обращения.
}

\NotesFrame{Слайд 12. Мобильное: создание обращения}{12}{
Это ключевой экран мобильного приложения. Житель выбирает категорию, описание, фотографирует проблему и указывает местоположение --- по GPS или вручную на карте.

\vspace{0.4em}

Если рядом уже есть похожее обращение, приложение предложит подтвердить его вместо создания дубля. Механизм дедупликации работает в радиусе 50~м.
}

\NotesFrame{Слайд 13. Мобильное: карта обращений}{13}{
На карте видны все обращения в виде маркеров. Цвет маркера показывает статус --- новое, в работе или решено.

\vspace{0.4em}

Можно фильтровать по категориям, найти себя на карте и быстро перейти к созданию нового обращения.
}

\NotesFrame{Слайд 14. Мобильное: список и карточка обращений}{14}{
Пользователь видит свои обращения в виде списка с вкладками --- все, активные и завершённые.

\vspace{0.4em}

При нажатии открывается полная карточка с фото, адресом, описанием и историей статусов. Другие жители могут нажать <<Тоже вижу проблему>>.
}

\NotesFrame{Слайд 15. Тестирование}{15}{
Проведено тестирование всех компонентов: автоматизированные тесты серверной части на xUnit, модульные тесты мобильного приложения на Vitest, тестирование API через Swagger~UI и функциональное тестирование веб-панели.

\vspace{0.4em}

Все тесты завершились успешно.
}

\NotesFrame{Слайд 16. Перспективы развития}{16}{
Дальнейшее развитие может включать push-уведомления на уровне ОС, офлайн-режим, расширение модуля городских событий, интеграцию с картографическими сервисами для обратного геокодирования и ролевую модель доступа для диспетчеров.
}

\NotesFrame{Слайд 17. Заключение}{17}{
В результате работы спроектирована и реализована программная система фиксации городских событий, включающая серверную часть, веб-панель диспетчера и мобильное приложение.

\vspace{0.4em}

Все задачи выполнены. Система обеспечивает полный цикл обработки обращений --- от подачи жителем до контроля диспетчером.

\vspace{0.4em}

Спасибо за внимание! Готова ответить на вопросы.
}

\end{document}
